\section{Rank, variability and shift calculations}
\label{app:theory}
These are standard calculations and implementation checks, not new distribution-free theorems. Condition on the training sample so that the score is fixed. With exchangeable continuous calibration and test scores, the test score has a uniform rank among $n+1$ scores. Inclusion in \eqref{eq:rank} is equivalent to rank at most $k$, proving $k/(n+1)$ coverage. Ties make the unrandomized comparison conservative. If $k>n$, the infinite threshold includes the score domain. Finite clr regions have the positive simplex as their domain. The implementation separately defines an infinite threshold as a whole-simplex fallback, including the boundary; this is an explicit rule, not evaluation of an undefined boundary clr score.

If scores are iid with continuous CDF $F$ and $k\le n$, the probability integral transform gives
\begin{equation}
 F(q)\sim\operatorname{Beta}(k,n+1-k),\qquad
 \operatorname{Var}\{F(q)\}=\frac{k(n+1-k)}{(n+1)^2(n+2)}.
\end{equation}
Thus even a perfectly marginally calibrated method varies across realized calibration samples. This beta statement needs iid sampling, stronger than general finite exchangeability. Given $q$, a test sample of size $m$ has independent Bernoulli indicators with success probability $F(q)$, so $\operatorname{Var}(\hat C\mid q)=F(q)(1-F(q))/m$. An independent finite test sample thus adds binomial noise. For calibration size $n$, test size $m$ and $\mu=k/(n+1)$, total variance of the test coverage is $v+(\mu(1-\mu)-v)/m$, where $v$ is the beta variance, conditional on the trained continuous score. This follows by total variance and $E[F(q)(1-F(q))]=\mu(1-\mu)-v$.

For coordinatewise Bonferroni calibration, the probability of any missed coordinate is at most $\sum_j\alpha/K=\alpha$. Naive per-coordinate error $\alpha$ only gives the lower bound $\max(0,1-K\alpha)$ on joint coverage. Dependence can improve this weak bound but does not supply the desired general guarantee.

A clr difference belongs to the sum-zero subspace. Any orthonormal basis of that subspace preserves its Euclidean norm, giving clr/ilr equivalence. Multiplicative perturbation followed by closure adds a fixed clr vector to both arguments and preserves their difference. These identities justify treating ilr and clr L2 as the same score geometry.

Let $c_0,c_1$ be group coverages of a fixed calibrated rule and $\pi=\Pr(G=1)$. Changing only the group mixture to $\pi'$ changes marginal coverage by
\begin{equation}
 C(\pi')-C(\pi)=(\pi'-\pi)(c_1-c_0).
 \label{eq:shift}
\end{equation}
This explains why hidden subgroup imbalance matters under this particular shift. The identity does not cover changes in within-group score laws. Mondrian rank guarantees transfer only because those laws are preserved.

The triangle simplex at $K=3$ has Euclidean area $\sqrt{3}/2$. A Euclidean ball centered at its barycenter with radius $r$ below its inradius cuts out a disk in its affine plane, giving area fraction $2\pi r^2/\sqrt{3}$. This supplies a known case for checking the uniform-simplex Monte Carlo estimator.
